\scanpage{25}
\mspage{\textcircled{7}}

\margin{$X_2=\{e_2,\unsure{\zeta|\xi})$}
Supposons \inserted{donc} $X_2$ de cardinal 2, \deleted{\unsure{donc}} et \deleted{\unsure{que}} $X_3=X'_3$ donc
\[
X_3=\{e_3\}\amalg\{d\unsure{\zeta|\xi},\rho d\unsure{\zeta|\xi},\rho^2d\unsure{\zeta|\xi}\},
\]
considérons les deux carrés

$(*)$
\diagram{\begin{tikzcd}[ampersand replacement=\&,column sep=large,row sep=large]
X_3\arrow[r,"{\delta_1=d=((x,y)\longmapsto(x,y,y))}"]\&X_2\arrow[d,dash,"\delta_0"]\\
X_2\arrow[u,"\delta_3"]\arrow[r,"\delta_0"]\&X_1=\{e_1\}
\end{tikzcd}}

et

$(**)$
\diagram{\begin{tikzcd}[ampersand replacement=\&,column sep=large,row sep=large]
X_2\&X_3\arrow[l,"p=p_{12}"']\\
\begin{gathered}X_1\\\{e_1\}=\end{gathered}\arrow[u,"{K_1=\delta_0\deleted{\unsure{K}}}"]\&X_2\arrow[l,"p_1"]\arrow[u,"{K'_2=((x,y)\longmapsto(x,x,y))}"']
\end{tikzcd}}
\inserted{$(=\delta_3)$}.

sont-ils cartésiens ? On a $\delta_3(x,y)=(x,x,y)\deleted{\unsure{=}}$,
\[
\deleted{d\sigma(x,y)}\quad d\sigma(x,y)=d(y,x)=(y,x,x)=\deleted{\unsure{\rho}}\rho^2\delta_3(x,y)
\]
i.e. $\delta_3=\rho^{\deleted{\unsure{2}}}d\sigma$. Donc en général (supposant seulement $\operatorname{card}X_1=1$ i.e. $X_1=\{e_1\}$), \deleted{\unsure{et quand}} \unsure{pour les}, quand a-t-on, \inserted{pour $\unsure{\zeta|\xi},\unsure{\zeta'|\xi'}\in X_2$,}
\[
d(\unsure{\zeta|\xi})=\rho^{\deleted{\unsure{2}}}d\underbrace{\sigma(\unsure{\zeta'|\xi'})}_{\unsure{\zeta''|\xi''}}\quad ?
\]
Il faut \unsure{s.} qu'on ait \deleted{\unsure{alors}}
$\unsure{\zeta|\xi}=\unsure{\zeta''|\xi''}=e_2$, \deleted{donc} $\boxed{\text{OK}}$.
\margin{d'où $\unsure{\zeta'|\xi'}=e_2$}

\deleted{$\unsure{\zeta|\xi}=\unsure{\zeta'|\xi'}=e_2$ \unsure{ce}} cas est bien cartésien, sans autre hypothèse. Le cas \unsure{$(**)$} s'écrit
\[
p_{\deleted{\unsure{12}}}^{-1}(e_2)\overset{?}=\delta_3(X_2)\quad[=\rho d\sigma(X_2)=\rho d(X_2)]
\]
\deleted{Mais $\unsure{p_{31}(x,y,z)}=\unsure{(z,x)}=\unsure{p\rho^2(x,y,z)}$ \unsure{i.e.} $\unsure{p_{31}=p\rho^2}$,}
\deleted{\unsure{puis}} i.e.
\[
p^{-1}(e_2)\overset{?}=\rho d(X_2)\qquad(\text{on a a priori }\supset)
\]
Si $X_3=X'_3\overset{\mathrm{def}}=\unsure{\mathfrak S_3|\mathcal G_3}\cdot d(X_2)$,
\deleted{\unsure{alors un objet de} $\unsure{X_3}$}
\inserted{\unsure{l'on est ramené à}}
\deleted{$X'_3$ \unsure{s'écrit}} prouvons que \deleted{\unsure{donc}} si $\unsure{\zeta|\xi}\in X'_3\setminus\{e_3\}$ satisfait à $p(\unsure{\zeta|\xi})=e_2$, on a $\unsure{\zeta|\xi}\in\rho d(X_2)$. Or $\unsure{\zeta|\xi}$ s'écrit canoniquement
\[
\unsure{\zeta|\xi}=\rho^i d(\eta),\quad\text{avec }\eta\in X_2^*,\quad\text{et }i\in\{0,1,2\},
\]
et on a
\[
p\unsure{\zeta|\xi}=\begin{cases}
\eta\ (\ne e_2)&\text{si }i=0\\
\sigma\eta\ (\ne e_2)&\text{si }\unsure{i=1|i=2}
\end{cases}
\]
\deleted{$\unsure{(p\sigma_3d)(x,y)}=\unsure{p\rho d(x,y)}=\unsure{\mathrm{id}}$}
donc $p\unsure{\zeta|\xi}=e_2\Longrightarrow i=1$, Qed.

\margin{24}
